\chapter{Your own commands, packages and themes}
\label{ch:macros}

\section{LaTeX as a language}\label{latex-as-a-language}

LaTeX is a programming language, and its commands behave like functions and variables.

\begin{lstlisting}[language=latex]
% A variable: one place to change a value
\newcommand{\studyarea}{Greater London}

% A function with one parameter, used as \term{surplus value}
\newcommand{\term}[1]{\textbf{#1}}

% A function with two parameters, where the first is optional
\newcommand{\money}[2][PKR]{#1~#2}   % \money{160} gives PKR 160
\end{lstlisting}

\texttt{\#1} is the first parameter. \texttt{\textbackslash{}renewcommand} changes an existing command.

\section{Environments and counters}\label{environments-and-counters}

\begin{lstlisting}[language=latex]
\newtheorem{example}{Example}[section]      % numbered within each section

\begin{example}[A coffee-shop worker]
  ...
\end{example}

\newcounter{step}
\stepcounter{step}\thestep
\end{lstlisting}

\section{Numbers from your analysis}\label{numbers-from-your-analysis}

Don\textquotesingle t type results by hand into a document. Have your analysis write a file of commands:

\begin{lstlisting}[language=latex]
% Written by a script. Do not edit by hand.
\newcommand{\meanDemand}{369}
\end{lstlisting}

Then \texttt{\textbackslash{}input\{numbers\}} and write \texttt{\textbackslash{}meanDemand} in the text. When the data changes, the document follows. This prevents the most common errors in scientific writing, as described in the \href{../../git-ci/}{Git and automation guide}.

\section{Your own style file}\label{your-own-style-file}

Put a shared look in a \texttt{.sty} file, then load it in any document.

\begin{lstlisting}[language=latex]
% mystyle.sty
\ProvidesPackage{mystyle}
\RequirePackage{xcolor}
\definecolor{brand}{HTML}{0D6A6E}
\end{lstlisting}

\begin{lstlisting}[language=latex]
\usepackage{mystyle}
\end{lstlisting}

For a whole document type, such as a thesis, write a \texttt{.cls} class file instead. Most journals and universities supply one. Keep your text separate from the style, so that moving to a different template only needs one change.

\section{Useful packages}\label{useful-packages}

{\def\LTcaptype{none} % do not increment counter
\begin{longtable}[]{@{}ll@{}}
\toprule\noalign{}
Package & Purpose \\
\midrule\noalign{}
\endhead
\bottomrule\noalign{}
\endlastfoot
\texttt{amsmath} & Mathematics \\
\texttt{geometry} & Page size and margins \\
\texttt{graphicx} & Images \\
\texttt{booktabs} & Clean tables \\
\texttt{longtable}, \texttt{tabularx} & Long and flexible tables \\
\texttt{xcolor} & Colour \\
\texttt{hyperref} & Links \\
\texttt{cleveref} & Automatic reference words \\
\texttt{natbib}, \texttt{biblatex} & Citations \\
\texttt{siunitx} & Numbers and units \\
\texttt{tikz}, \texttt{pgfplots} & Diagrams and plots \\
\texttt{glossaries} & Glossaries and acronyms \\
\texttt{tcolorbox} & Coloured boxes \\
\texttt{beamer} & Slides, as a class \\
\end{longtable}
}

Load only what you use, and load \texttt{hyperref} late, as it changes many others.

